% ======================================================================
% CAPÍTULO 6 — RESULTADOS E DISCUSSÃO
% ======================================================================
\chapter{RESULTADOS E DISCUSSÃO}

Este capítulo apresenta os resultados da validação do sistema, obtidos pela
execução das suítes de testes automatizados descritas na metodologia, e discute
esses resultados à luz dos requisitos e dos critérios de aceitação definidos. Os
valores reportados correspondem a execuções reais realizadas no ambiente de
desenvolvimento do projeto.

\section{Estratégia de validação executada}

A validação foi organizada em três frentes, conforme a camada do sistema:

\begin{itemize}
  \item \textbf{Backend (Python):} testes unitários (FSM, PID, rampa,
  persistência, enlace serial), testes da API REST e do WebSocket e testes
  funcionais de ciclo completo e de segurança, executados com o \texttt{pytest};
  \item \textbf{Firmware (\daq):} testes \emph{host-based} da biblioteca
  portável \texttt{daqcore} (parser JSON, watchdog e pulso da bomba), executados
  sem hardware no ambiente \texttt{native} do PlatformIO;
  \item \textbf{IHM Web:} testes unitários dos \emph{stores} de telemetria, de
  acionamento manual e de tema, executados com o Vitest.
\end{itemize}

A Tabela~\ref{tab:testes} consolida os resultados das execuções.

\begin{table}[htbp]
\centering
\caption{Resultado da execução das suítes de testes automatizados.}
\label{tab:testes}
\begin{tabular}{lcccc}
\toprule
\textbf{Camada} & \textbf{Ferramenta} & \textbf{Arquivos} & \textbf{Testes} & \textbf{Resultado} \\
\midrule
Backend (Python) & pytest & 8 & 40 & 40 aprovados \\
Firmware (\daq) & PlatformIO (host) & 1 & 14 & 14 aprovados \\
IHM Web & Vitest & 4 & 13 & 13 aprovados \\
\midrule
\textbf{Total} & --- & 13 & 67 & 67 aprovados \\
\bottomrule
\end{tabular}
\end{table}

\par\vspace{-2pt}{\footnotesize Fonte: execu\c{c}\~oes realizadas em 05/09/2026 no ambiente de desenvolvimento do projeto.}

\section{Testes unit\'arios e de integra\c{c}\~ao}

\subsection{Backend}

Os testes do backend cobrem os principais m\'odulos da l\'ogica de supervis\~ao:

\begin{itemize}
  \item \textbf{FSM:} valida\c{c}\~ao das transi\c{c}\~oes entre estados, da
  matriz de acionamento em cada fase e do comportamento dos eventos de parada e
  de emerg\^encia;
  \item \textbf{PID e rampa:} valida\c{c}\~ao do c\'alculo da sa\'ida do
  controlador, da satura\c{c}\~ao, do anti-windup, da taxa de aquecimento e da
  sele\c{c}\~ao das duas regi\~oes de controle da rampa (raz\~ao de taxas abaixo
  de $0\,\celsius$ e PID acima);
  \item \textbf{Persist\^encia:} valida\c{c}\~ao da escrita e leitura at\^omicas,
  da restaura\c{c}\~ao do backup e da valida\c{c}\~ao de faixas;
  \item \textbf{Enlace serial:} valida\c{c}\~ao do codec JSON, da
  serializa\c{c}\~ao/desserializa\c{c}\~ao e do tratamento de pacotes
  malformados;
  \item \textbf{API e integra\c{c}\~ao:} valida\c{c}\~ao dos endpoints de
  configura\c{c}\~ao e de comando e do fluxo de ciclo completo com simulador de
  \daq via portas seriais virtuais.
\end{itemize}

\subsection{Firmware}

Os testes do firmware (\texttt{test\_daqcore}) validam: o parser de pacotes JSON
(escrita e handshake de configura\c{c}\~ao), a rejei\c{c}\~ao de JSON inv\'alido,
a satura\c{c}\~ao do PWM, a constru\c{c}\~ao do pacote de resposta, o
comportamento do watchdog (n\~ao disparo antes do tempo limite, disparo ap\'os o
tempo limite e rearme por nova batida) e a l\'ogica de pulso da bomba (pulso nas
transi\c{c}\~oes liga/desliga e encerramento ap\'os o tempo definido).

\subsection{IHM Web}

Os testes do frontend validam a l\'ogica de assinatura e consumo da telemetria
WebSocket, o sincronismo do estado dos controles manuais com a telemetria e a
persist\^encia da prefer\^encia de tema, garantindo que os componentes reajam
corretamente aos dados recebidos.

\section{Testes funcionais do ciclo completo}

Os testes funcionais exercitam o ciclo autom\'atico $T_0 \to T_3$ com o
simulador do \daq, verificando que a FSM percorre as quatro fases na ordem
correta e que a matriz de acionamento \'e aplicada conforme o esperado em cada
fase (Tabela~\ref{tab:matriz}). A valida\c{c}\~ao confirmou que: (i) o ciclo
somente inicia a partir do estado seguro; (ii) as dura\c{c}\~oes das fases
seguem os par\^ametros configurados; e (iii) a telemetria publicada ao longo do
ciclo cont\'em o estado, o progresso da etapa e o progresso geral do ciclo,
habilitando o Diagrama de Tempos e a barra de progresso da IHM.

Adicionalmente, o roteiro de estresse de comunica\c{c}\~ao a 4~Hz, executado com
o simulador via \texttt{socat}, exercita o enlace serial de forma cont\'inua,
verificando a integridade da troca de pacotes e a estabilidade do loop de
controle ao longo do tempo.

\section{Testes de seguran\c{c}a}

Os cen\'arios de seguran\c{c}a simulam a perda de comunica\c{c}\~ao serial e
verificam as duas camadas de prote\c{c}\~ao:

\begin{itemize}
  \item \textbf{Watchdog do firmware:} quando nenhum pacote v\'alido \'e recebido
  dentro do tempo limite (1~s), o firmware comuta os atuadores para o estado
  seguro, desligando os fornos;
  \item \textbf{Watchdog do backend:} quando o \daq deixa de responder por mais
  de 1~s em uma fase ativa, o backend dispara o evento de emerg\^encia e a FSM
  retorna ao estado seguro.
\end{itemize}

Os testes confirmam que, em ambos os casos, os fornos s\~ao desligados dentro do
intervalo de seguran\c{c}a especificado ($\le 1$~s), atendendo ao crit\'erio de
aceita\c{c}\~ao de seguran\c{c}a.

\section{Discuss\~ao}

\subsection{Atendimento aos objetivos de valida\c{c}\~ao}

Os resultados obtidos indicam que a implementa\c{c}\~ao atende aos objetivos
funcionais e de seguran\c{c}a definidos na metodologia, com 67 testes aprovados
distribu\'idos pelas tr\^es camadas do sistema. Tr\^es observa\c{c}\~oes merecem
destaque:

\begin{enumerate}[label=\arabic*)]
  \item \textbf{Corre\c{c}\~ao da l\'ogica de controle:} os testes de FSM, PID e
  rampa validam matematicamente a estrat\'egia de controle mista, incluindo o
  tratamento da satura\c{c}\~ao e a sele\c{c}\~ao da regi\~ao de opera\c{c}\~ao,
  o que \'e condi\c{c}\~ao necess\'aria --- ainda que n\~ao suficiente --- para a
  qualidade da rampa em bancada;
  \item \textbf{Seguran\c{c}a por constru\c{c}\~ao:} a valida\c{c}\~ao dos
  watchdogs em ambas as camadas demonstra que o requisito de parada segura
  autom\'atica \'e atendido no n\'ivel de software, independentemente do
  hardware;
  \item \textbf{Integra\c{c}\~ao entre camadas:} os testes de ciclo completo com
  simulador confirmam a coer\^encia entre a FSM, o protocolo serial e a
  telemetria consumida pela IHM, reduzindo o risco de falhas de integra\c{c}\~ao
  quando o hardware real for conectado.
\end{enumerate}

\subsection{Compara\c{c}\~ao com os crit\'erios de aceita\c{c}\~ao}

Dos cinco crit\'erios de aceita\c{c}\~ao consolidados na
Tabela~\ref{tab:criterios}, tr\^es puderam ser verificados no ambiente de
valida\c{c}\~ao (disponibilidade da comunica\c{c}\~ao, persist\^encia e
seguran\c{c}a do watchdog) e dois dependem de medi\c{c}\~ao em bancada
(linearidade da rampa e estabilidade do Forno 2). \'E importante destacar que
esses \'ultimos s\~ao crit\'erios de \emph{desempenho t\'ermico} do processo
f\'isico e n\~ao podem ser atestados por testes de software: eles exigem a
calibra\c{c}\~ao da curva de aquecimento do Tubo U e a verifica\c{c}\~ao
experimental das curvas de temperatura, etapas de continuidade deste trabalho.

\subsection{Limita\c{c}\~oes}

A principal limita\c{c}\~ao da valida\c{c}\~ao realizada \'e a aus\^encia de
medi\c{c}\~oes em bancada com o processo f\'isico. Os testes automatizados
validam a correta execu\c{c}\~ao da l\'ogica e dos mecanismos de seguran\c{c}a,
mas n\~ao substituem a verifica\c{c}\~ao experimental da din\^amica t\'ermica.
Al\'em disso, a raz\~ao de taxas empregada na regi\~ao de malha aberta depende da
calibra\c{c}\~ao da taxa de aquecimento do sistema (curva ``Taxa $\times$ PWM''),
que deve ser medida no equipamento real; na aus\^encia desse dado, a
implementa\c{c}\~ao adota uma aproxima\c{c}\~ao que precisa ser confirmada.
Registra-se, ainda, a indefini\c{c}\~ao quanto ao circuito integrado exato de
convers\~ao do termopar (MAX6675 ou MAX31855), que n\~ao altera a l\'ogica de
controle, mas deve ser confirmada na documenta\c{c}\~ao f\'isica do equipamento.

\subsection{Implica\c{c}\~oes}

No plano t\'ecnico, os resultados demonstram a viabilidade da arquitetura
proposta para substituir o sistema legado, com vantagens de custo, flexibilidade
e seguran\c{c}a. No plano operacional, a persist\^encia dos par\^ametros e o
registro de tend\^encias criam condi\c{c}\~oes para a repetibilidade e a
auditoria do m\'etodo, requisitos valorizados em ambiente anal\'itico. A etapa de
valida\c{c}\~ao em bancada \'e o passo natural para confirmar os crit\'erios de
desempenho t\'ermico e para calibrar a rampa do Tubo U.
